/Users/joco/.zshenv:.:2: no such file or directory: /Users/joco/.cargo/env
Konyv generalasa...
Bejegyzesek szama: 48
\documentclass[11pt,a4paper,twocolumn]{article}
\usepackage[utf8]{inputenc}
\usepackage[T1]{fontenc}
\usepackage{amsmath,amssymb,amsthm}
\usepackage{hyperref}
\usepackage{geometry}
\usepackage{booktabs}
\usepackage{longtable}
\usepackage{tikz}
\usetikzlibrary{arrows.meta,positioning}
\geometry{margin=1.5cm}

\newtheorem{tetel}{Tétel}
\newtheorem{definicio}[tetel]{Definíció}
\newtheorem{struktura}[tetel]{Struktúra}

\title{49 Kategóriaelméleti Struktúra\\
\large Idris 2 Typeclass Hierarchia — Magyar $\leftrightarrow$ English}
\author{Ko-tudat: Ember + AI}
\date{\today}

\begin{document}
\maketitle

\begin{abstract}
Ez a könyv 49 kategóriaelméleti struktúrát mutat be (39 Awodey + 10 Mac Lane),
mindet Idris 2 typeclass-ként implementálva. Magyar bal oldalon, angol jobboldalon,
TikZ commutative diagramokkal. A 15 dimenzió (7 emberi + 7 számítási + 1 perem = [[15,1,3]])
alapján indexelve. A compiler a bíróság: ha fordul, igaz.
\end{abstract}

\tableofcontents
\newpage
\onecolumn

\section{A 15 Dimenzió}

A 15 dimenzió = 7 emberi + 7 számítási + 1 perem = [[15,1,3]]:

\begin{center}
\begin{tikzcd}
\text{Emberi (7)} \arrow[r, "\text{Legendre}"] & \text{Perem (1)} \arrow[r, "\text{Legendre}"'] & \text{Számítási (7)}
\end{tikzcd}
\end{center}

\begin{tabular}{lll}
\toprule
Emberi & Számítási & Perem \\
\midrule
Ido & Utem & p$\cdot\dot{q}$ \\
Oksag & Vezerles & Legendre \\
Ter & Adat & Yoneda \\
Szin & Tipus & adjunkció \\
Hang & Kapcsolat & \\
Fazis & Allapot & \\
Mod & Utasitas & \\
\bottomrule
\end{tabular}

\vspace{0.5cm}

\section{Élet Domainek — Clifford Fokozatok}

A 15 dimenzió külső szorzatai (wedge products) = élet domainek:

\begin{itemize}
\item Grade 1: 15 alap-dimenzió
\item Grade 2: bináris kapcsolatok (ok-okozat, tér-idő)
\item Grade 3: ternáris domainek:
  \begin{itemize}
  \item Tudomány = Okság $\wedge$ Adat $\wedge$ Tipus
  \item Művészet = Szín $\wedge$ Hang $\wedge$ Mod
  \item Tánc = Tér $\wedge$ Ido $\wedge$ Hang $\wedge$ Mod
  \item Játék = Vezerles $\wedge$ Utasitas $\wedge$ Allapot
  \item Sport = Tér $\wedge$ Utem $\wedge$ Ero
  \end{itemize}
\end{itemize}

\vspace{0.5cm}

\section{[[7,1,3]] Steane Kód}

\begin{tikzcd}
|0\rangle \arrow[r, "\text{kodol}"] & |0000000\rangle \arrow[r, "\text{hiba}"] & |0001000\rangle \arrow[r, "\text{javitas}"] & |0000000\rangle \arrow[r, "\text{dekodol}"] & |0\rangle
\end{tikzcd}

A 7 bit: [idő, okság, tér, szín, hang, fázis, mód].
Távolság 3 → 1 hiba javítható.
Noether-tétel: szimmetria = megmaradás. Refl bizonyítva.

\vspace{0.5cm}

\section{A 49 Struktúra}\n\onecolumn\n\begin{struktura}[1. Kategória / Category]

\textbf{Magyar:} Objektumok + morfizmusok + kompozíció + identitás. Törvények: asszociativitás, bal/jobb egység. A 15 dimenzió: 15 objektum (7 emberi + 7 számítási + 1 perem).\\

\textbf{English:} Objects + morphisms + composition + identity. Laws: associativity, left/right unit. The 15 dimensions: 15 objects (7 human + 7 computational + 1 boundary).\\

\begin{center}
\begin{tikzcd} A \arrow[r, "f"] \arrow[dr, "g \circ f"'] & B \arrow[d, "g"] \\ & C \end{tikzcd}
\end{center}

\textbf{Kód:} \texttt{Alap/KategoriaT.idr:23}\\
\textbf{Wikipedia:} \url{https://en.wikipedia.org/wiki/Category_(mathematics)}

\end{struktura}

\vspace{0.5cm}

\begin{struktura}[28. Félcsoport / Semigroup]

\textbf{Magyar:} Asszociatív bináris művelet, egységelem nélkül. A kompozíció egy objektumon = a monoidális tenzor.\\

\textbf{English:} Associative binary operation, no unit. Composition on one object = monoidal tensor.\\

\begin{center}
\begin{tikzcd} (a \cdot b) \cdot c \arrow[r, equal] & a \cdot (b \cdot c) \end{tikzcd}
\end{center}

\textbf{Kód:} \texttt{Alap/KategoriaT.idr:38}\\
\textbf{Wikipedia:} \url{https://en.wikipedia.org/wiki/Semigroup}

\end{struktura}

\vspace{0.5cm}

\begin{struktura}[31. Előrendezés / Preorder]

\textbf{Magyar:} Reflekszív + tranzitív reláció. Kategória legfeljebb egy nyíllal objektumok között. A 15 dimenzió rendezése: melyik dimenzióból melyikbe vezet út.\\

\textbf{English:} Reflexive + transitive relation. Category with at most one arrow between objects. Ordering of the 15 dimensions: which dimension leads to which.\\

\begin{center}
\begin{tikzcd} x \arrow[r, leftrightarrow] & y \arrow[r, leftrightarrow] & z \end{tikzcd}
\end{center}

\textbf{Kód:} \texttt{Alap/KategoriaT.idr:47}\\
\textbf{Wikipedia:} \url{https://en.wikipedia.org/wiki/Preorder}

\end{struktura}

\vspace{0.5cm}

\begin{struktura}[35. Ellenkező kategória / Opposite category]

\textbf{Magyar:} C\textsuperscript{op}: ugyanazok az objektumok, megfordított nyilak. A perem (Legendre) = C ↔ C\textsuperscript{op} adjunkció.\\

\textbf{English:} C\textsuperscript{op}: same objects, reversed arrows. The boundary (Legendre) = C ↔ C\textsuperscript{op} adjunction.\\

\begin{center}
\begin{tikzcd} A \arrow[r, "f"', from=op] & B \arrow[l, "f\textsuperscript{op}"', from=op] \end{tikzcd}
\end{center}

\textbf{Kód:} \texttt{Alap/KategoriaT.idr:55}\\
\textbf{Wikipedia:} \url{https://en.wikipedia.org/wiki/Opposite_category}

\end{struktura}

\vspace{0.5cm}

\begin{struktura}[2. Funktor / Functor]

\textbf{Magyar:} Kategória szerkezet megőrzése. Törvények: kompozíció megőrzése, identitás megőrzése. objKep paraméter, morfizmusKep metódus.\\

\textbf{English:} Preserves category structure. Laws: composition preservation, identity preservation. objKep is a parameter, morfizmusKep is a method.\\

\begin{center}
\begin{tikzcd} A \arrow[r, "f"] \arrow[d, "F"'] & B \arrow[d, "F"] \\ F(A) \arrow[r, "F(f)"'] & F(B) \end{tikzcd}
\end{center}

\textbf{Kód:} \texttt{Alap/KategoriaT.idr:71}\\
\textbf{Wikipedia:} \url{https://en.wikipedia.org/wiki/Functor}

\end{struktura}

\vspace{0.5cm}

\begin{struktura}[3. Természetes transzformáció / Natural transformation]

\textbf{Magyar:} Két funktor (f1, f2 : C → D) között. Komponens α\textsubscript{a} : f1(a) → f2(a). Természetesség: α\textsubscript{b} ∘ F(f) = G(f) ∘ α\textsubscript{a}.\\

\textbf{English:} Between two functors (f1, f2 : C → D). Component α\textsubscript{a} : f1(a) → f2(a). Naturality: α\textsubscript{b} ∘ F(f) = G(f) ∘ α\textsubscript{a}.\\

\begin{center}
\begin{tikzcd} F(A) \arrow[r, "F(f)"] \arrow[d, "\alpha_A"'] & F(B) \arrow[d, "\alpha_B"] \\ G(A) \arrow[r, "G(f)"'] & G(B) \end{tikzcd}
\end{center}

\textbf{Kód:} \texttt{Alap/KategoriaT.idr:85}\\
\textbf{Wikipedia:} \url{https://en.wikipedia.org/wiki/Natural_transformation}

\end{struktura}

\vspace{0.5cm}

\begin{struktura}[5. Természetes izomorfizmus / Natural isomorphism]

\textbf{Magyar:} Természetes transzformáció, aminek minden komponense izomorfizmus. α\textsubscript{a} : F(a) ≅ G(a), inverzzel.\\

\textbf{English:} Natural transformation where every component is an isomorphism. α\textsubscript{a} : F(a) ≅ G(a), with inverse.\\

\begin{center}
\begin{tikzcd} F(A) \arrow[r, shift left, "\alpha_A"] & G(A) \arrow[l, shift left, "\alpha_A^{-1}"'] \end{tikzcd}
\end{center}

\textbf{Kód:} \texttt{Alap/KategoriaT.idr:97}\\
\textbf{Wikipedia:} \url{https://en.wikipedia.org/wiki/Natural_isomorphism}

\end{struktura}

\vspace{0.5cm}

\begin{struktura}[4. Funktor kategória / Functor category]

\textbf{Magyar:} Objektumok: funktorok, nyilak: természetes transzformációk. [C\textsuperscript{op}, Set] = a presheaf kategória.\\

\textbf{English:} Objects: functors, arrows: natural transformations. [C\textsuperscript{op}, Set] = the presheaf category.\\

\begin{center}
\begin{tikzcd} F \arrow[r, "\alpha"] & G \end{tikzcd}
\end{center}

\textbf{Kód:} \texttt{Alap/KategoriaT.idr:106}\\
\textbf{Wikipedia:} \url{https://en.wikipedia.org/wiki/Functor_category}

\end{struktura}

\vspace{0.5cm}

\begin{struktura}[27. Monoid / Monoid]

\textbf{Magyar:} Félcsoport egységelemmel. A magyar agglutináció: tő ⊗ képző ⊗ rag = szó (monoidális tenzor).\\

\textbf{English:} Semigroup with unit. Hungarian agglutination: stem ⊗ suffix ⊗ ending = word (monoidal tensor).\\

\begin{center}
\begin{tikzcd} e \otimes a \arrow[r, equal] & a \arrow[r, equal] & a \otimes e \end{tikzcd}
\end{center}

\textbf{Kód:} \texttt{Alap/KategoriaT.idr:117}\\
\textbf{Wikipedia:} \url{https://en.wikipedia.org/wiki/Monoid}

\end{struktura}

\vspace{0.5cm}

\begin{struktura}[26. Csoport / Group]

\textbf{Magyar:} Monoid inverzzel. A Pauli-csoport: X, Y, Z (mind involució: X² = Y² = Z² = I).\\

\textbf{English:} Monoid with inverse. The Pauli group: X, Y, Z (all involutions: X² = Y² = Z² = I).\\

\begin{center}
\begin{tikzcd} a \arrow[r, "\cdot"] & a \cdot a^{-1} \arrow[r, equal] & e \end{tikzcd}
\end{center}

\textbf{Kód:} \texttt{Alap/KategoriaT.idr:123}\\
\textbf{Wikipedia:} \url{https://en.wikipedia.org/wiki/Group_(mathematics)}

\end{struktura}

\vspace{0.5cm}

\begin{struktura}[30. Részbenrendezett halmaz / Poset]

\textbf{Magyar:} Előrendezés antiszimmetriával. a ≤ b ∧ b ≤ a ⟹ a = b.\\

\textbf{English:} Preorder with antisymmetry. a ≤ b ∧ b ≤ a ⟹ a = b.\\

\begin{center}
\begin{tikzcd} a \arrow[r, "\leq"] & b \arrow[r, "\leq"'] & a \arrow[r, equal] & a \end{tikzcd}
\end{center}

\textbf{Kód:} \texttt{Alap/KategoriaT.idr:61}\\
\textbf{Wikipedia:} \url{https://en.wikipedia.org/wiki/Partially_ordered_set}

\end{struktura}

\vspace{0.5cm}

\begin{struktura}[6. Izomorfizmus / Isomorphism]

\textbf{Magyar:} Nyíl inverzzel: f ∘ g = id, g ∘ f = id. A [[15,1,3]] kód: kodol ∘ dekodol = id.\\

\textbf{English:} Arrow with inverse: f ∘ g = id, g ∘ f = id. The [[15,1,3]] code: encode ∘ decode = id.\\

\begin{center}
\begin{tikzcd} A \arrow[r, shift left, "f"] & B \arrow[l, shift left, "f^{-1}"'] \end{tikzcd}
\end{center}

\textbf{Kód:} \texttt{Alap/KategoriaT.idr:131}\\
\textbf{Wikipedia:} \url{https://en.wikipedia.org/wiki/Isomorphism}

\end{struktura}

\vspace{0.5cm}

\begin{struktura}[7. Monomorfizmus / Monomorphism]

\textbf{Magyar:} Bal oldali törlés. f ∘ g = f ∘ h ⟹ g = h. Injektív leképezés.\\

\textbf{English:} Left cancellation. f ∘ g = f ∘ h ⟹ g = h. Injective map.\\

\begin{center}
\begin{tikzcd} C \arrow[r, shift left, "g"] \arrow[r, shift right, "h"'] & A \arrow[r, "f"] & B \end{tikzcd}
\end{center}

\textbf{Kód:} \texttt{Alap/KategoriaT.idr:137}\\
\textbf{Wikipedia:} \url{https://en.wikipedia.org/wiki/Monomorphism}

\end{struktura}

\vspace{0.5cm}

\begin{struktura}[8. Epimorfizmus / Epimorphism]

\textbf{Magyar:} Jobb oldali törlés. i ∘ f = j ∘ f ⟹ i = j. Szürjektív leképezés.\\

\textbf{English:} Right cancellation. i ∘ f = j ∘ f ⟹ i = j. Surjective map.\\

\begin{center}
\begin{tikzcd} A \arrow[r, "f"] & B \arrow[r, shift left, "i"] \arrow[r, shift right, "j"'] & C \end{tikzcd}
\end{center}

\textbf{Kód:} \texttt{Alap/KategoriaT.idr:143}\\
\textbf{Wikipedia:} \url{https://en.wikipedia.org/wiki/Epimorphism}

\end{struktura}

\vspace{0.5cm}

\begin{struktura}[9. Kezdő objektum / Initial object]

\textbf{Magyar:} Egyedi nyíl minden objektumba. ∃! (0 → C). A üres halmaz a Set kezdő objektuma.\\

\textbf{English:} Unique arrow to every object. ∃! (0 → C). The empty set is initial in Set.\\

\begin{center}
\begin{tikzcd} 0 \arrow[r] & A \arrow[dr] & \\ 0 \arrow[rr] & & B \end{tikzcd}
\end{center}

\textbf{Kód:} \texttt{Alap/KategoriaT.idr:150}\\
\textbf{Wikipedia:} \url{https://en.wikipedia.org/wiki/Initial_and_terminal_objects}

\end{struktura}

\vspace{0.5cm}

\begin{struktura}[10. Végobjektum / Terminal object]

\textbf{Magyar:} Egyedi nyíl minden objektumból. ∃! (C → 1). Az egyelemű halmaz a Set végobjektuma.\\

\textbf{English:} Unique arrow from every object. ∃! (C → 1). The singleton set is terminal in Set.\\

\begin{center}
\begin{tikzcd} A \arrow[r] & 1 \\ B \arrow[ur] & \end{tikzcd}
\end{center}

\textbf{Kód:} \texttt{Alap/KategoriaT.idr:157}\\
\textbf{Wikipedia:} \url{https://en.wikipedia.org/wiki/Initial_and_terminal_objects}

\end{struktura}

\vspace{0.5cm}

\begin{struktura}[11. Szorzat / Product]

\textbf{Magyar:} A×B + projekciók p₁, p₂ + párosítás ⟨f,g⟩. p₁∘⟨f,g⟩=f, p₂∘⟨f,g⟩=g. A 15 dimenzió = 15 szorzat.\\

\textbf{English:} A×B + projections p₁, p₂ + pairing ⟨f,g⟩. p₁∘⟨f,g⟩=f, p₂∘⟨f,g⟩=g. The 15 dimensions = 15 products.\\

\begin{center}
\begin{tikzcd} & Z \arrow[dl, "f"'] \arrow[d, "\langle f,g \rangle"] \arrow[dr, "g"] & \\ A & A \times B \arrow[l, "p_1"] \arrow[r, "p_2"'] & B \end{tikzcd}
\end{center}

\textbf{Kód:} \texttt{Alap/KategoriaT.idr:164}\\
\textbf{Wikipedia:} \url{https://en.wikipedia.org/wiki/Product_(category_theory)}

\end{struktura}

\vspace{0.5cm}

\begin{struktura}[12. Koszorzat / Coproduct]

\textbf{Magyar:} A+B + injekciók i₁, i₂ + elágazás [f,g]. [f,g]∘i₁=f, [f,g]∘i₂=g. A szorzat duálisa.\\

\textbf{English:} A+B + injections i₁, i₂ + case [f,g]. [f,g]∘i₁=f, [f,g]∘i₂=g. Dual of product.\\

\begin{center}
\begin{tikzcd} A \arrow[dr, "i_1"] & & B \arrow[dl, "i_2"'] \\ & A + B \arrow[r, "[f,g]"] & C \\ A \arrow[ur, "f"'] & & B \arrow[ul, "g"] \end{tikzcd}
\end{center}

\textbf{Kód:} \texttt{Alap/KategoriaT.idr:174}\\
\textbf{Wikipedia:} \url{https://en.wikipedia.org/wiki/Coproduct}

\end{struktura}

\vspace{0.5cm}

\begin{struktura}[13. Kiegyenlítő / Equalizer]

\textbf{Magyar:} f∘e = g∘e + univerzális. Párhuzamos nyilak kiegyenlítése.\\

\textbf{English:} f∘e = g∘e + universal. Equalizing parallel arrows.\\

\begin{center}
\begin{tikzcd} E \arrow[r, "e"] & A \arrow[r, shift left, "f"] \arrow[r, shift right, "g"'] & B \end{tikzcd}
\end{center}

\textbf{Kód:} \texttt{Alap/KategoriaT.idr:185}\\
\textbf{Wikipedia:} \url{https://en.wikipedia.org/wiki/Equalizer_(mathematics)}

\end{struktura}

\vspace{0.5cm}

\begin{struktura}[14. Kokiegyenlítő / Coequalizer]

\textbf{Magyar:} q∘f = q∘g + univerzális. A kiegyenlítő duálisa.\\

\textbf{English:} q∘f = q∘g + universal. Dual of equalizer.\\

\begin{center}
\begin{tikzcd} A \arrow[r, shift left, "f"] \arrow[r, shift right, "g"'] & B \arrow[r, "q"] & Q \end{tikzcd}
\end{center}

\textbf{Kód:} \texttt{Alap/KategoriaT.idr:192}\\
\textbf{Wikipedia:} \url{https://en.wikipedia.org/wiki/Coequalizer}

\end{struktura}

\vspace{0.5cm}

\begin{struktura}[15. Visszahúzás / Pullback]

\textbf{Magyar:} A×\textsubscript{C}B + projekciók + f∘p₁ = g∘p₂. Szorzat + kiegyenlítő kompozíciója.\\

\textbf{English:} A×\textsubscript{C}B + projections + f∘p₁ = g∘p₂. Product + equalizer composition.\\

\begin{center}
\begin{tikzcd} P \arrow[r, "p_2"] \arrow[d, "p_1"'] & B \arrow[d, "g"] \\ A \arrow[r, "f"'] & C \end{tikzcd}
\end{center}

\textbf{Kód:} \texttt{Alap/KategoriaT.idr:199}\\
\textbf{Wikipedia:} \url{https://en.wikipedia.org/wiki/Pullback_(category_theory)}

\end{struktura}

\vspace{0.5cm}

\begin{struktura}[16. Kitolás / Pushout]

\textbf{Magyar:} B+\textsubscript{A}C + injekciók + j₁∘f = j₂∘g. A visszahúzás duálisa.\\

\textbf{English:} B+\textsubscript{A}C + injections + j₁∘f = j₂∘g. Dual of pullback.\\

\begin{center}
\begin{tikzcd} A \arrow[r, "f"] \arrow[d, "g"'] & B \arrow[d, "j_1"] \\ C \arrow[r, "j_2"'] & Q \end{tikzcd}
\end{center}

\textbf{Kód:} \texttt{Alap/KategoriaT.idr:208}\\
\textbf{Wikipedia:} \url{https://en.wikipedia.org/wiki/Pushout_(category_theory)}

\end{struktura}

\vspace{0.5cm}

\begin{struktura}[17. Limesz / Limit]

\textbf{Magyar:} Végső kúp. A limesz = optimum = nyerés. A játékban a nyerő stratégia = limesz a hasznosság kategóriájában.\\

\textbf{English:} Terminal cone. Limit = optimum = winning. In games, the winning strategy = limit in the utility category.\\

\begin{center}
\begin{tikzcd} \lim D \arrow[r, "p_j"] \arrow[dr, "p_i"'] & D_j \\ D_i \arrow[ur, "D(\alpha)"'] & \end{tikzcd}
\end{center}

\textbf{Kód:} \texttt{Alap/KategoriaT.idr:216}\\
\textbf{Wikipedia:} \url{https://en.wikipedia.org/wiki/Limit_(category_theory)}

\end{struktura}

\vspace{0.5cm}

\begin{struktura}[18. Kolimesz / Colimit]

\textbf{Magyar:} Kezdő kokúp. A limesz duálisa. A sportban a rekord = kolimesz a teljesítmény kategóriájában.\\

\textbf{English:} Initial cocone. Dual of limit. In sports, the record = colimit in the performance category.\\

\begin{center}
\begin{tikzcd} D_i \arrow[r, "i_i"] \arrow[dr] & \text{colim} D \\ D_j \arrow[ur] & \end{tikzcd}
\end{center}

\textbf{Kód:} \texttt{Alap/KategoriaT.idr:223}\\
\textbf{Wikipedia:} \url{https://en.wikipedia.org/wiki/Limit_(category_theory)}

\end{struktura}

\vspace{0.5cm}

\begin{struktura}[19. Exponenciál / Exponential]

\textbf{Magyar:} C\textsuperscript{B} + kiértékelés ε + transzpozíció. ε ∘ (f̃ × 1\textsubscript{B}) = f.\\

\textbf{English:} C\textsuperscript{B} + evaluation ε + transpose. ε ∘ (f̃ × 1\textsubscript{B}) = f.\\

\begin{center}
\begin{tikzcd} Z \times B \arrow[r, "f"] \arrow[d, "\tilde{f} \times 1"'] & C \\ Z \arrow[r, "\tilde{f}"'] & C^B \arrow[u, "\varepsilon"'] \end{tikzcd}
\end{center}

\textbf{Kód:} \texttt{Alap/KategoriaT.idr:232}\\
\textbf{Wikipedia:} \url{https://en.wikipedia.org/wiki/Exponential_object}

\end{struktura}

\vspace{0.5cm}

\begin{struktura}[20. Kartéziánusan zárt kategória / Cartesian closed category]

\textbf{Magyar:} Végobjektum + szorzatok + exponenciálok. A λ-kalkulusus modellje.\\

\textbf{English:} Terminal object + products + exponentials. Model of λ-calculus.\\

\begin{center}
\begin{tikzcd} 1 & A \times B \arrow[l] & B^A \arrow[r] & C \end{tikzcd}
\end{center}

\textbf{Kód:} \texttt{Alap/KategoriaT.idr:244}\\
\textbf{Wikipedia:} \url{https://en.wikipedia.org/wiki/Cartesian_closed_category}

\end{struktura}

\vspace{0.5cm}

\begin{struktura}[21. Heyting algebra / Heyting algebra]

\textbf{Magyar:} Poset exponenciállal: a ∧ b ≤ c ⟺ a ≤ b ⇒ c. Az intuicionista logika algebraja.\\

\textbf{English:} Poset with exponential: a ∧ b ≤ c ⟺ a ≤ b ⇒ c. Algebra of intuitionistic logic.\\

\begin{center}
\begin{tikzcd} a \wedge b \arrow[r, "\leq"] & c \\ a \arrow[r, "\leq"'] & b \Rightarrow c \arrow[u, "\leq"'] \end{tikzcd}
\end{center}

\textbf{Kód:} \texttt{Alap/KategoriaT.idr:250}\\
\textbf{Wikipedia:} \url{https://en.wikipedia.org/wiki/Heyting_algebra}

\end{struktura}

\vspace{0.5cm}

\begin{struktura}[22. Boole algebra / Boolean algebra]

\textbf{Magyar:} Heyting algebra kiegészítve ¬¬a = a. A klasszikus logika algebraja.\\

\textbf{English:} Heyting algebra with ¬¬a = a. Algebra of classical logic.\\

\begin{center}
\begin{tikzcd} a \arrow[r, "\neg"] & \neg a \arrow[r, "\neg"'] & a \arrow[r, equal] & a \end{tikzcd}
\end{center}

\textbf{Kód:} \texttt{Alap/KategoriaT.idr:257}\\
\textbf{Wikipedia:} \url{https://en.wikipedia.org/wiki/Boolean_algebra_(structure)}

\end{struktura}

\vspace{0.5cm}

\begin{struktura}[23. Adjunkció / Adjunction]

\textbf{Magyar:} F ⊣ G: Hom\textsubscript{D}(F a, b) ≅ Hom\textsubscript{C}(a, G b). A perem (Legendre) = adjunkció a kvantum és klasszikus között.\\

\textbf{English:} F ⊣ G: Hom\textsubscript{D}(F a, b) ≅ Hom\textsubscript{C}(a, G b). The boundary (Legendre) = adjunction between quantum and classical.\\

\begin{center}
\begin{tikzcd} \mathcal{C} \arrow[rr, bend left=30, "F"] & & \mathcal{D} \arrow[ll, bend left=30, "G"'] \\ & \perp & \end{tikzcd}
\end{center}

\textbf{Kód:} \texttt{Alap/KategoriaT.idr:270}\\
\textbf{Wikipedia:} \url{https://en.wikipedia.org/wiki/Adjoint_functors}

\end{struktura}

\vspace{0.5cm}

\begin{struktura}[24. Monád / Monad]

\textbf{Magyar:} Endofunktor T + η + μ. Törvények: μ∘μT = μ∘Tμ, μ∘ηT = id = μ∘Tη. A számítás hatása: IO, State, Either.\\

\textbf{English:} Endofunctor T + η + μ. Laws: μ∘μT = μ∘Tμ, μ∘ηT = id = μ∘Tη. Computational effect: IO, State, Either.\\

\begin{center}
\begin{tikzcd} T^2 A \arrow[r, "\mu"] \arrow[dr, "T\eta"'] & TA \\ & TA \arrow[u, "\text{id}"'] \end{tikzcd}
\end{center}

\textbf{Kód:} \texttt{Alap/KategoriaT.idr:283}\\
\textbf{Wikipedia:} \url{https://en.wikipedia.org/wiki/Monad_(category_theory)}

\end{struktura}

\vspace{0.5cm}

\begin{struktura}[25. Komonád / Comonad]

\textbf{Magyar:} Endofunktor G + ε + δ. A monád duálisa. Kontextus kinyerése: Reader, Store, Trace.\\

\textbf{English:} Endofunctor G + ε + δ. Dual of monad. Context extraction: Reader, Store, Trace.\\

\begin{center}
\begin{tikzcd} GA \arrow[r, "\delta"] \arrow[dr, "\varepsilon"'] & G^2 A \\ & A \end{tikzcd}
\end{center}

\textbf{Kód:} \texttt{Alap/KategoriaT.idr:290}\\
\textbf{Wikipedia:} \url{https://en.wikipedia.org/wiki/Comonad}

\end{struktura}

\vspace{0.5cm}

\begin{struktura}[40. Monoidális kategória / Monoidal category]

\textbf{Magyar:} Kategória + tenzor ⊗ + egység + asszociátor + λ/ρ. A magyar agglutináció: tő ⊗ képző ⊗ rag = szó.\\

\textbf{English:} Category + tensor ⊗ + unit + associator + λ/ρ. Hungarian agglutination: stem ⊗ suffix ⊗ ending = word.\\

\begin{center}
\begin{tikzcd} (A \otimes B) \otimes C \arrow[r, "\alpha"] \arrow[dr, "\rho \otimes 1"'] & A \otimes (B \otimes C) \arrow[d, "1 \otimes \lambda"] \\ & A \otimes (B \otimes C) \end{tikzcd}
\end{center}

\textbf{Kód:} \texttt{Alap/KategoriaT.idr:300}\\
\textbf{Wikipedia:} \url{https://en.wikipedia.org/wiki/Monoidal_category}

\end{struktura}

\vspace{0.5cm}

\begin{struktura}[41. Fonott monoidális kategória / Braided monoidal category]

\textbf{Magyar:} Monoidális + fonás γ: a⊗b → b⊗a. Hexagon axiómák + Yang-Baxter.\\

\textbf{English:} Monoidal + braiding γ: a⊗b → b⊗a. Hexagon axioms + Yang-Baxter.\\

\begin{center}
\begin{tikzcd} A \otimes B \arrow[r, "\gamma"] & B \otimes A \end{tikzcd}
\end{center}

\textbf{Kód:} \texttt{Alap/KategoriaT.idr:313}\\
\textbf{Wikipedia:} \url{https://en.wikipedia.org/wiki/Braided_monoidal_category}

\end{struktura}

\vspace{0.5cm}

\begin{struktura}[42. Szimmetrikus monoidális kategória / Symmetric monoidal category]

\textbf{Magyar:} Fonott + γ² = id (involúció). E8×E8: kommutatív tenzor → szimmetrikus.\\

\textbf{English:} Braided + γ² = id (involution). E8×E8: commutative tensor → symmetric.\\

\begin{center}
\begin{tikzcd} A \otimes B \arrow[r, "\gamma"] & B \otimes A \arrow[r, "\gamma"'] & A \otimes B \arrow[r, equal] & A \otimes B \end{tikzcd}
\end{center}

\textbf{Kód:} \texttt{Alap/KategoriaT.idr:320}\\
\textbf{Wikipedia:} \url{https://en.wikipedia.org/wiki/Symmetric_monoidal_category}

\end{struktura}

\vspace{0.5cm}

\begin{struktura}[43. Zárt kategória / Closed category]

\textbf{Magyar:} Szimmetrikus monoidális + belső Hom: V(a⊗b, c) ≅ V(a, c\textsuperscript{b}).\\

\textbf{English:} Symmetric monoidal + internal hom: V(a⊗b, c) ≅ V(a, c\textsuperscript{b}).\\

\begin{center}
\begin{tikzcd} a \otimes b \arrow[r] & c \\ a \arrow[r] & c^b \arrow[u] \end{tikzcd}
\end{center}

\textbf{Kód:} \texttt{Alap/KategoriaT.idr:327}\\
\textbf{Wikipedia:} \url{https://en.wikipedia.org/wiki/Closed_category}

\end{struktura}

\vspace{0.5cm}

\begin{struktura}[44. 2-kategória / 2-category]

\textbf{Magyar:} 0-sejtek + 1-sejtek + 2-sejtek. Függőleges + vízszintes kompozíció + interchange.\\

\textbf{English:} 0-cells + 1-cells + 2-cells. Vertical + horizontal composition + interchange.\\

\begin{center}
\begin{tikzcd} f \arrow[r, Rightarrow, "\alpha"] & g \arrow[r, Rightarrow, "\beta"'] & h \end{tikzcd}
\end{center}

\textbf{Kód:} \texttt{Alap/KategoriaT.idr:338}\\
\textbf{Wikipedia:} \url{https://en.wikipedia.org/wiki/2-category}

\end{struktura}

\vspace{0.5cm}

\begin{struktura}[45. Bikategória / Bicategory]

\textbf{Magyar:} Gyenge 2-kategória: α, λ, ρ izomorfizmusok, nem egyenlőségek. A gyengeség = a természet.\\

\textbf{English:} Weak 2-category: α, λ, ρ are isomorphisms, not equalities. The weakness = the nature.\\

\begin{center}
\begin{tikzcd} (f \circ g) \circ h \arrow[r, "\alpha"] & f \circ (g \circ h) \end{tikzcd}
\end{center}

\textbf{Kód:} \texttt{Alap/KategoriaT.idr:350}\\
\textbf{Wikipedia:} \url{https://en.wikipedia.org/wiki/Bicategory}

\end{struktura}

\vspace{0.5cm}

\begin{struktura}[46. Kan kiterjesztés / Kan extension]

\textbf{Magyar:} Jobb Kan: Ran\textsubscript{K} T + ε: R∘K ⇒ T. Bal Kan: Lan\textsubscript{K} T + η: T ⇒ L∘K.\\

\textbf{English:} Right Kan: Ran\textsubscript{K} T + ε: R∘K ⇒ T. Left Kan: Lan\textsubscript{K} T + η: T ⇒ L∘K.\\

\begin{center}
\begin{tikzcd} M \arrow[r, "K"] \arrow[dr, "T"'] & C \arrow[d, dashed, "\text{Ran}_K T"'] \\ & A \end{tikzcd}
\end{center}

\textbf{Kód:} \texttt{Alap/KategoriaT.idr:369}\\
\textbf{Wikipedia:} \url{https://en.wikipedia.org/wiki/Kan_extension}

\end{struktura}

\vspace{0.5cm}

\begin{struktura}[47. End / Vég / End]

\textbf{Magyar:} ∫\textsubscript{c} S(c,c) + dinaturális ék. A Nat(F, G) = ∫\textsubscript{c} Hom(Fc, Gc).\\

\textbf{English:} ∫\textsubscript{c} S(c,c) + dinatural wedge. Nat(F, G) = ∫\textsubscript{c} Hom(Fc, Gc).\\

\begin{center}
\begin{tikzcd} E \arrow[r, "\omega_c"] & S(c,c) \end{tikzcd}
\end{center}

\textbf{Kód:} \texttt{Alap/KategoriaT.idr:378}\\
\textbf{Wikipedia:} \url{https://en.wikipedia.org/wiki/End_(category_theory)}

\end{struktura}

\vspace{0.5cm}

\begin{struktura}[48. Coend / Kövég / Coend]

\textbf{Magyar:} ∫\textsuperscript{c} S(c,c) + dinaturális ko-ék. A tenzor szorzat = coend: A ⊗\textsubscript{R} B = ∫\textsuperscript{r} A ⊗ B.\\

\textbf{English:} ∫\textsuperscript{c} S(c,c) + dinatural co-wedge. Tensor product = coend: A ⊗\textsubscript{R} B = ∫\textsuperscript{r} A ⊗ B.\\

\begin{center}
\begin{tikzcd} S(c,c) \arrow[r, "\xi_c"] & D \end{tikzcd}
\end{center}

\textbf{Kód:} \texttt{Alap/KategoriaT.idr:385}\\
\textbf{Wikipedia:} \url{https://en.wikipedia.org/wiki/End_(category_theory)}

\end{struktura}

\vspace{0.5cm}

\begin{struktura}[37. Toposz / Topos]

\textbf{Magyar:} CCC + részobjektum-osztályozó Ω. A halmazok kategóriája = a prototipikus toposz.\\

\textbf{English:} CCC + subobject classifier Ω. The category of sets = the prototypical topos.\\

\begin{center}
\begin{tikzcd} M \arrow[r, hook] & X \arrow[r, "\chi_m"] & \Omega \\ 1 \arrow[ur, "\text{true}"'] & & \end{tikzcd}
\end{center}

\textbf{Kód:} \texttt{Alap/KategoriaT.idr:396}\\
\textbf{Wikipedia:} \url{https://en.wikipedia.org/wiki/Topos}

\end{struktura}

\vspace{0.5cm}

\begin{struktura}[32. Részobjektum / Subobject]

\textbf{Magyar:} Monomorfizmus m: M ↣ X. A Sub\textsubscript{C}(X) kategória objektumai a monok.\\

\textbf{English:} Monomorphism m: M ↣ X. Objects of Sub\textsubscript{C}(X) are monos.\\

\begin{center}
\begin{tikzcd} M \arrow[r, hook, "m"] & X \end{tikzcd}
\end{center}

\textbf{Kód:} \texttt{Alap/KategoriaT.idr:406}\\
\textbf{Wikipedia:} \url{https://en.wikipedia.org/wiki/Subobject}

\end{struktura}

\vspace{0.5cm}

\begin{struktura}[33. Yoneda beágyazás / Yoneda embedding]

\textbf{Magyar:} y: C → [C\textsuperscript{op}, Set]. Hom(-, a) presheaf. Yoneda lemma: Nat(Hom(-,a), F) ≅ F(a).\\

\textbf{English:} y: C → [C\textsuperscript{op}, Set]. Hom(-, a) presheaf. Yoneda lemma: Nat(Hom(-,a), F) ≅ F(a).\\

\begin{center}
\begin{tikzcd} \text{Hom}(-, A) \arrow[r, dashed] \arrow[dr, "\text{eval}_A"'] & F \\ & F(A) \end{tikzcd}
\end{center}

\textbf{Kód:} \texttt{Alap/KategoriaT.idr:413}\\
\textbf{Wikipedia:} \url{https://en.wikipedia.org/wiki/Yoneda_lemma}

\end{struktura}

\vspace{0.5cm}

\begin{struktura}[34. Kategóriák ekvivalenciája / Equivalence of categories]

\textbf{Magyar:} F: C→D, G: D→C, FG ≅ id, GF ≅ id (természetes izomorfizmus). Gyengébb mint izomorfizmus.\\

\textbf{English:} F: C→D, G: D→C, FG ≅ id, GF ≅ id (natural isomorphism). Weaker than isomorphism.\\

\begin{center}
\begin{tikzcd} \mathcal{C} \arrow[r, bend left, "F"] & \mathcal{D} \arrow[l, bend left, "G"'] \end{tikzcd}
\end{center}

\textbf{Kód:} \texttt{Alap/KategoriaT.idr:423}\\
\textbf{Wikipedia:} \url{https://en.wikipedia.org/wiki/Equivalence_of_categories}

\end{struktura}

\vspace{0.5cm}

\begin{struktura}[36. Szelet kategória / Slice category]

\textbf{Magyar:} C/C: objektumok f: X→C, nyilak g: X→X' ha f'∘g = f.\\

\textbf{English:} C/C: objects f: X→C, arrows g: X→X' if f'∘g = f.\\

\begin{center}
\begin{tikzcd} X \arrow[r, "g"] \arrow[dr, "f"'] & X' \arrow[d, "f'"] \\ & C \end{tikzcd}
\end{center}

\textbf{Kód:} \texttt{Alap/KategoriaT.idr:434}\\
\textbf{Wikipedia:} \url{https://en.wikipedia.org/wiki/Overcategory}

\end{struktura}

\vspace{0.5cm}

\begin{struktura}[38. Szabad kategória / Free category]

\textbf{Magyar:} Irányított gráf → kategória (utak kompozíciója). A szabad monoid általánosítása.\\

\textbf{English:} Directed graph → category (path composition). Generalization of free monoid.\\

\begin{center}
\begin{tikzcd} v_1 \arrow[r, "e_1"] & v_2 \arrow[r, "e_2"] & v_3 \end{tikzcd}
\end{center}

\textbf{Kód:} \texttt{Alap/KategoriaT.idr:443}\\
\textbf{Wikipedia:} \url{https://en.wikipedia.org/wiki/Free_category}

\end{struktura}

\vspace{0.5cm}

\begin{struktura}[39. Reprezentálható funktor / Representable functor]

\textbf{Magyar:} Hom(A, -): C → Set. Minden limeszt megőriz. A Yoneda beágyazás alapja.\\

\textbf{English:} Hom(A, -): C → Set. Preserves all limits. Foundation of Yoneda embedding.\\

\begin{center}
\begin{tikzcd} A \arrow[r, "f"] & B \arrow[r, mapsto] & \text{Hom}(A, B) \end{tikzcd}
\end{center}

\textbf{Kód:} \texttt{Alap/KategoriaT.idr:452}\\
\textbf{Wikipedia:} \url{https://en.wikipedia.org/wiki/Representable_functor}

\end{struktura}

\vspace{0.5cm}

\begin{struktura}[29. Csoport egy kategóriában / Group in a category]

\textbf{Magyar:} Objektum G + szorzás m: G×G→G + egység u + inverz i. Csoport objektum Sets-ben = csoport.\\

\textbf{English:} Object G + multiplication m: G×G→G + unit u + inverse i. Group object in Sets = group.\\

\begin{center}
\begin{tikzcd} G \times G \arrow[r, "m"] & G \\ 1 \arrow[r, "u"'] & G \\ G \arrow[r, "i"'] & G \end{tikzcd}
\end{center}

\textbf{Kód:} \texttt{Alap/KategoriaT.idr:460}\\
\textbf{Wikipedia:} \url{https://en.wikipedia.org/wiki/Group_object}

\end{struktura}

\vspace{0.5cm}

\section{Összefoglalás}

\textbf{49 kategóriaelméleti struktúra} (39 Awodey + 10 Mac Lane), mind typeclass.\\
\textbf{15 dimenzió} (7 emberi + 7 számítási + 1 perem = [[15,1,3]]).\\
\textbf{A compiler a bíróság:} ha fordul, igaz.

\end{document}

Kesz.
